\section*{Unterstützungsleistungen}
\addcontentsline{toc}{section}{Unterstützungsleistungen}

\begin{itemize}
    \item \textbf{Anthropic, Claude, KI-Sprachmodell:} Art der Unterstützung: Code-Review, Fehlersuche und Debugging sowie Unterstützung beim Entwerfen, Formulieren und kritischen Überarbeiten von Text und Software. Vorschläge und Ausgaben wurden vom Verfasser geprüft und bei Bedarf verändert.
    \item \textbf{OpenAI, KI-gestützte Werkzeuge (einschließlich ChatGPT/Codex):} Art der Unterstützung: Code-Review, Fehlersuche und Debugging sowie Unterstützung beim Entwerfen, Formulieren und kritischen Überarbeiten von Text und Software. Vorschläge und Ausgaben wurden vom Verfasser geprüft und bei Bedarf verändert.
    \item \textbf{TODO: [Vor- und Nachname der betreuenden Lehrkraft], Betreuungslehrkraft, Gymnasium Mülheim an der Ruhr, Mülheim an der Ruhr:} Art der Unterstützung: Fachliche Begleitung und Beratung bei Themenfindung und Strukturierung der schriftlichen Arbeit. \textbf{Vor Abgabe vervollständigen.}
    \item \textbf{Sparkassenstiftung Mülheim an der Ruhr, Mülheim an der Ruhr:} Stiftung; Art der Unterstützung: Sponsoring-Anträge für mechatronische Hardwarekomponenten eingereicht.
\end{itemize}

Die Datenerfassung, Ausführung der Trainingsläufe und Bewertung der Resultate lagen beim Verfasser. KI-Werkzeuge dienten als unterstützende Arbeitsmittel; die Verantwortung für Auswahl, Prüfung und Darstellung der Inhalte verbleibt beim Verfasser. Die Unterstützungsleistungen sind zusätzlich vollständig in der Jufo-Wettbewerbsverwaltung einzutragen.
